\begin{frame}{자주 하는 실수 5가지}
  각각이 Week 2--7의 어떤 원칙을 깨는지 짚어두면, 8.2절에서
  바로 반박 가능한 질문으로 바꿀 수 있다:
  \begin{enumerate}
    \item \kb{테스트로 모델 고르기} (``살짝 훔쳐보기'') --- 선택
      편향을 test에 그대로 전이. \textit{리뷰 질문}: ``val F1
      순위와 test F1 순위가 같은가요? 시드 3개로 확인할 수
      있나요?''
    \item \kb{정규화/결측치 통계를 전체 데이터로 fit} --- 6.3절의
      ``정규화 통계량 누수''. \texttt{fit(X\_train)} $\to$
      \texttt{transform(X\_val)}, \texttt{transform(X\_test)}
      가 정석.
    \item \kb{베이스라인 없이 성능 보고} --- 클래스 비율 90:10에서
      ``정확도 91\%''는 아무것도 말하지 않는다.
    \item \kb{시계열/그룹 데이터의 무작위 분할} --- ``비현실적으로
      좋은'' 성능, \textbf{배포(새 환자/미래 시간)에서는 완전히
      붕괴}.
    \item \kb{학습 곡선 없이 하이퍼파라미터 보고} --- ``lr=0.1을
      썼다''가 왜 0.05·0.3이 아니라 0.1인지 val 곡선 없이는
      확인 불가(M3가 학습 곡선을 요구하는 이유).
  \end{enumerate}
\end{frame}
